% TODO make hw stylesheet
\documentclass[11pt]{scrartcl}
\usepackage[a4paper, total={6in, 8in}]{geometry}
\usepackage[usenames,dvipsnames,svgnames]{xcolor}
\usepackage[shortlabels]{enumitem}
\usepackage[framemethod=TikZ]{mdframed}
\usepackage{amsmath,amssymb,amsthm}
\usepackage[colorlinks]{hyperref}
\usepackage{multicol}
\usepackage{tabularx}

\hypersetup{
	colorlinks=true,
	linkcolor=blue,
	filecolor=magenta,      
	urlcolor=magenta,
	pdftitle={MAT 320 HW}
}

\usepackage{microtype}
\usepackage{mathtools}
\usepackage[headsepline]{scrlayer-scrpage}
\usepackage{thmtools}
\usepackage{todonotes}

\mdfdefinestyle{mdgreenbox}{linecolor=ForestGreen,backgroundcolor=ForestGreen!5,
	linewidth=2pt,rightline=false,leftline=true,topline=false,bottomline=false,}
\declaretheoremstyle[headfont=\bfseries\sffamily\color{ForestGreen!70!black},
mdframed={style=mdgreenbox},headpunct={ --- },]{thmgreenbox}

\declaretheorem[style=thmgreenbox,name=Claim,numbered=no]{claim*}
\declaretheorem[style=thmgreenbox,name=Lemma,numbered=no]{lemma*}

\addtolength{\textheight}{3.14cm}
\providecommand{\ol}{\overline}
\providecommand{\ul}{\underline}
\providecommand{\eps}{\varepsilon}
\providecommand{\half}{\frac{1}{2}}
\providecommand{\dang}{\measuredangle} %% Directed angle
\providecommand{\CC}{\mathbb C}
\providecommand{\FF}{\mathbb F}
\providecommand{\NN}{\mathbb N}
\providecommand{\QQ}{\mathbb Q}
\providecommand{\RR}{\mathbb R}
\providecommand{\ZZ}{\mathbb Z}
\providecommand{\dg}{^\circ}
\providecommand{\ii}{\item}
\providecommand{\setdiff}{\smallsetminus}
\providecommand{\alert}[1]{{\sffamily\textbf{\textcolor{blue}{#1}}}}

\reversemarginpar

\theoremstyle{problem}
\newtheorem{innercustomthm}{Problem}
\newenvironment{problem}[1]
{\renewcommand\theinnercustomthm{#1}\innercustomthm}
{\endinnercustomthm}

\theoremstyle{remark}
\newtheorem{remarknumbered}{Remark}
\newenvironment{remark}
{\renewcommand\theremarknumbered{\hspace{-\the\fontdimen2\font\space}}\remarknumbered}
{\endremarknumbered}

\newenvironment{soln}{\begin{proof}[Solution]}{\end{proof}}
\newenvironment{walkthrough}{\noindent \textbf{Walkthrough.}}{\medskip}
\newlist{walk}{enumerate}{3}
\setlist[walk]{label=\bfseries (\alph*)}

\providecommand{\pow}{\operatorname{Pow}}
\providecommand{\vocab}[1]{{\textbf{\textcolor{ForestGreen}{#1}}}}
\providecommand{\lcm}{\operatorname{lcm}}
\providecommand{\abs}[1]{\left|#1\right|}

\usepackage{asymptote}
\begin{asydef}
	defaultpen(fontsize(10pt));
	unitsize(1cm);
	size(8cm); // set a reasonable default
	usepackage("amsmath");
	usepackage("amssymb");
	settings.tex="pdflatex";
	settings.outformat="pdf";
	// Replacement for olympiad+cse5 which is not standard
	import geometry;
	// recalibrate fill and filldraw for conics
	void filldraw(picture pic = currentpicture, conic g, pen fillpen=defaultpen, pen drawpen=defaultpen)
	{ filldraw(pic, (path) g, fillpen, drawpen); }
	void fill(picture pic = currentpicture, conic g, pen p=defaultpen)
	{ filldraw(pic, (path) g, p); }
	// some geometry
	pair foot(pair P, pair A, pair B) { return foot(triangle(A,B,P).VC); }
	pair orthocenter(pair A, pair B, pair C) { return orthocentercenter(A,B,C); }
	pair centroid(pair A, pair B, pair C) { return (A+B+C)/3; }
	// cse5 abbreviations
	path CP(pair P, pair A) { return circle(P, abs(A-P)); }
	path CR(pair P, real r) { return circle(P, r); }
	pair IP(path p, path q) { return intersectionpoints(p,q)[0]; }
	pair OP(path p, path q) { return intersectionpoints(p,q)[1]; }
	path Line(pair A, pair B, real a=0.6, real b=a) { return (a*(A-B)+A)--(b*(B-A)+B); }
	// cse5 more useful functions
	picture CC() {
		picture p=rotate(0)*currentpicture;
		currentpicture.erase();
		return p;
	}
	pair MP(Label s, pair A, pair B = plain.S, pen p = defaultpen) {
		Label L = s;
		L.s = "$"+s.s+"$";
		label(L, A, B, p);
		return A;
	}
	pair Drawing(Label s = "", pair A, pair B = plain.S, pen p = defaultpen) {
		dot(MP(s, A, B, p), p);
		return A;
	}
	path Drawing(path g, pen p = defaultpen, arrowbar ar = None) {
		draw(g, p, ar);
		return g;
	}
\end{asydef}

\usepackage{mathrsfs}

\ihead{\footnotesize MAT 320 HW06}
\ohead{\footnotesize October 16, 2024}

\title{\vspace{-2em}MAT 320 HW06}
\author{\large Aditya Pahuja}
\date{\large Due 10/16}
\begin{document}
\maketitle

\begin{problem}{1}
	Prove that $x = \cos(x)$ for some $x\in(0, \pi/2)$.
\end{problem}

\begin{soln}
	Consider the function $f(x) = \cos(x) - x$. Assuming that $\cos(x)$
	and $x$ are both continuous on $[0, \pi/2]$, we can see that
	$f(x)$ is also continuous on $[0, \pi/2]$. Now, since
	\[f(0) = \cos(0) - 0 = 1 > 0,
	\qquad f(\pi/2) = \cos(\pi/2) - \pi/2 = -\pi/2 < 0,\]
	the intermediate value theorem says that there exists an
	$x_0\in[0, \pi/2]$ such that $f(x_0) = 0$.
	Clearly $f(0)$ and $f(\pi/2)$ are not zero,
	so $x_0 \ne 0$ and $x_0\ne\pi/2$, thus $x_0\in(0, \pi/2)$, which means that
	we have found a value of $x$ (namely $x_0$) in the open interval
	$(0, \pi/2)$ such that
	\[f(x) = \cos(x) - x = 0 \iff \cos(x) = x,\]
	as desired.
\end{soln}

\begin{problem}{2}
	Show that a polynomial $P(x)$ of odd degree
	has at least one zero.
\end{problem}

\begin{soln}
	Let $P(x) = a_0 + a_1x + a_2x^2 + \cdots + a_nx^n$
	where $n$ is an odd integer and the $a_i$ are real numbers
	such that $a_n\ne 0$.
	
	\begin{claim*}
		There exist real numbers $r_+$, $r_-$
		such that $P(r_+) > 0$ and $P(r_-) < 0$.
	\end{claim*}
	\begin{proof}
		Consider the expression
		\[\frac{P(x)}{x^n} = a_n + \frac{a_{n-1}}{x}
		+ \cdots + \frac{a_1}{x^{n-1}} + \frac{a_0}{x^n}
		= a_n +  \sum_{k=0}^{n-1} \frac{a_k}{x^{n-k}}.\]
		Since
		\[\lim_{x\to\infty} \frac{c}{x^k} = 0\]
		for any real $c$ and positive integer $k$, we can say that
		\[\lim_{x\to\infty} \frac{P(x)}{x^n} = a_n
		+ \lim_{x\to\infty}  \sum_{k=0}^{n-1} \frac{a_k}{x^{n-k}} = a_n,\]
		as the terms with $x$ in the denominator
		grow arbitrarily small when $x$ increases.
		Similarly, we have
		\[\lim_{x\to\infty} \frac{P(-x)}{x^n} =
		\lim_{x\to\infty} -a_n + \sum_{k=0}^{n-1} \frac{a_k\cdot(-1)^k}{x^{n-k}}
		= -a_n,
		\]
		where we use the fact that $n$ is odd to conclude that
		$a_n(-x)^n = -a_nx^n$.
		
		Now, suppose without loss of generality
		that $a_n > 0$ and let $\eps = a_n$.
		By the definition of limit, the first limit tells us that
		there exists an $N_1\in\NN$ such that
		$x > N_1$ implies $\abs{\frac{P(x)}{x^n} - a_n} < \eps$.
		Similarly, the second limit posits the existence of an $N_2\in\NN$
		such that $x > N_2$ implies that $\abs{\frac{P(-x)}{x^n} + a_n} < \eps$.
		Equivalently, the inequalities rewrite as
		\[0 = a_n - \eps < \frac{P(x)}{x^n} < a_n + \eps < 2a_n,\qquad
		-2a_n = -a_n - \eps < \frac{P(-x)}{x^n} < -a_n + \eps = 0.\]
		Since $x^n$ is positive whenever $x > \max(N_1, N_2)$,
		we can say that $P(x)$ is positive and $P(-x)$ is negative
		for any choice of $x > \max(N_1, N_2)$.
		Then, we can take $r_+ = x$ and $r_- = -x$ to get the claim.
	\end{proof}
	
	To finish, since $P(r_+)$ is positive and $P(r_-)$ is negative,
	we can use the fact that polynomials are continuous
	to conclude that, by the intermediate value theorem,
	there is an $r_0$ between $r_-$ and $r_+$
	for which $P(r_0) = 0$.
\end{soln}

\begin{problem}{3}
	Prove that a continuous function $f:\RR\to\RR$
	that takes on no value more than twice
	takes some value exactly once.
\end{problem}

\begin{soln}
	Let $a$ and $b$ be real numbers such that $a < b$ and
	$f(a) = f(b) = L$; if no such $a$ and $b$ exist,
	then $f$ only takes each value at most once,
	so we are immediately done by picking any value
	in the image of $f$.
	By the Weierstrass theorem, there exists $u$ and $v$ such that,
	for all $x\in[a, b]$,
	\[f(u) \le f(x) \le f(v).\]
	Now, $f(u)$ and $f(v)$ cannot both be equal to $L$,
	as otherwise we would get $f(x) = L$ for all $x\in[a, b]$,
	and there are more than two values of $x$ in $[a, b]$,
	which contradicts our assumption about $f$.
	So, suppose WLOG that $f(u) < L$ (as opposed to $f(v) > L$).
	
	I claim that $f$ takes on the value $m = f(u)$ exactly once.
	Suppose for the sake of contradiction that $f(x) = m$
	has two solutions, say $u$ and another real number $w$.
	Assume WLOG $u < w$ and consider two cases:
	\begin{itemize}
		\ii $w\in[a, b]$. In this case, we know that
		$f(x) \ge m$ for all $x\in[a, b]$, so in particular
		$f(x) \ge m$ for all $x\in[u, w]$.
		The Weierstrass theorem tells us that, more specifically,
		$f$ attains some maximum $M$ on $[u, w]$; that is,
		there exists a real $t\in[u, w]$ such that $m \le f(x) \le f(t) = M$
		when $x\in[u,w]$.
		
		Now, let $r$ be a real number in the interval $(m, \min(L, M))$.
		I will find three values $x_1$, $x_2$, and $x_3$ such that
		$f(x_1) = f(x_2) = f(x_3) = r$.
		First, observe that
		\[f(u) = f(w) < r < f(a) = f(b),\]
		so, by the intermediate value theorem on $f$ over the intervals $[a, u]$
		and $[w, b]$, there exists $x_1\in[a, u]$ and $x_2\in[w, b]$
		such that $f(x_1) = f(x_2) = r$.
		Note that $r\ne m$ means that $x_1$ and $x_2$ are distinct
		from $u$ and $w$, so really we can say $x_1 \in [a, u)$
		and $x_2\in (w, b]$.
		For the third value, we see that, by the intermediate value theorem,
		there exists $x_3$ within the interval $[u, t]$
		such that $f(x_3) = r\in[m, M] = [f(u), f(t)]$.
		
		This means $r$ is achieved by $f$ more than twice,
		which is a contradiction.
		
		\ii $w\notin[a, b]$. Our assumption that $w > u$ means that
		$w > b$. Pick an $r$ in the open interval $(m, L)$.
		Then, by intermediate value theorem, the following values exist:
		\begin{itemize}
			\ii[$\circ$] $x_1\in[a, u]$ satisfying $f(x_1) = r$, since
			$r\in[f(u), f(a)] = [m, L]$.
			\ii[$\circ$] $x_2\in[u, b]$ satisfying $f(x_2) = r$, since
			$r\in[f(u), f(b)] = [m, L]$.
			\ii[$\circ$] $x_3\in[b, w]$ satisfying $f(x_3) = r$, since
			$r\in[f(w), f(b)] = [m, L]$.
		\end{itemize}
		Moreover, no two of the $x_i$ can be equal, as
		$x_1 = x_2 = u$ would imply $r = m$ and $x_2 = x_3$ would imply
		$r = L$, both of which contradict $r\in(m, L)$
		(and obviously $x_1 < x_2 < x_3$ so $x_1\ne x_3$).
		Thus, we again have found an $r$ that is achieved by $f$
		more than twice, which is a contradiction.
	\end{itemize}

	To conclude, $f(x) = m$ must only have one real solution $u$,
	as desired.
\end{soln}

\begin{problem}{4}
	Let $d$ and $d'$ be two equivalent metrics on $X$. Show that:
	\begin{itemize}
		\ii a set $A\subseteq X$ is open in $(X, d)$ if and only if
		it is open in $(X, d')$.
		\ii a set $A\subseteq X$ is closed in $(X, d)$ if and only if
		it is closed in $(X, d')$.
		\ii a function $f\colon X\to\RR$ is continuous with respect to
		the metric $d$ if and only if it is continuous with respect to
		the metric $d'$.
		\ii a sequence $(x_n)$ of points in $X$ converges to some $\ell\in X$
		with respect to the metric $d$ if and only if
		it converges with respect to $\ell$ with respect to $d'$.
	\end{itemize}
\end{problem}

\begin{soln}
	We can rephrase the equivalence of $d$ and $d'$ as saying
	that there exist positive reals $a$ and $b$ such that
	\[a \le \frac{d(x, y)}{d'(x, y)} \le b\]
	for all $x,\,y\in X$. That is, the ratio of the distances
	between $x$ and $y$ in the two metrics is bounded within
	an interval with positive real endpoints.
	This lets us see that equivalence of metrics is symmetric ---
	if $\frac{d}{d'}$ is in $[a, b]$,
	then $\frac{d'}{d}$ is always in $[\frac 1b, \frac1a]$.
	Thus, it suffices to prove just the forward direction
	(assuming $d$ has the property) for each of the statements,
	as the argument in the backward direction will be identical.
	
	\paragraph{Openness.}
	If $A\subseteq X$ is open with respect to $d$, then, for each $x\in A$,
	there exists an $\eps$ such that $B(x,\eps_x)\subseteq A$.
	That is, for each $y\in X$,
	\[d(x, y) < \eps_x \implies y \in A.\]
	Using the equivalence of $d$ and $d'$,
	\[a\cdot d'(x, y)\le d(x, y) < \eps_x\implies y \in A.\]
	This can be rephrased as
	\[d'(x, y) < \frac{\eps_x}{a}\implies y \in A,\]
	meaning that, for each $x\in A$, we have found an open ball
	$B(x, \frac{\eps_x}{a})$ (defined in the metric space $(X, d')$)
	centered at $x$ that is a subset of $A$.
	Thus, $A$ is also open with respect to $d'$
	(and as mentioned above, the converse holds
	because we can just swap $d$ and $d'$ in the proof).
	
	\paragraph{Closedness.}
	If $A\subseteq X$ is closed with respect to $d$,
	then we can equivalently say that its complement
	$X\smallsetminus A$ is open with respect to $d$.
	By the previous section, this is equivalent to
	$X\smallsetminus A$ is being open with respect to $d'$,
	which is equivalent to its complement $A$
	being closed with respect to $d'$.
	
	\paragraph{Continuity.}
	Suppose $f\colon X\to \RR$ is continuous in $(X, d)$.
	Then, for each $c\in X$ and $\eps > 0$, there exists $\delta$ such that
	\[d(x, c) < \delta \implies d_{\RR}(f(x), f(c)) = |f(x) - f(c)| < \eps\]
	(where I am assuming that $\RR$ is equipped with the standard
	absolute value metric $d_\RR(x, y) = |x - y|$).
	The condition $d(x, c) < \delta$ implies that
	\[a\cdot d'(x, c) \le d(x, c) < \delta \implies
	d'(x, c) < \frac{\delta}{a}.\]
	That is, given $\eps$ and $c$, if $\delta$ is the quantity mentioned above,
	then we can say that there exists $\delta' = \frac{\delta}{a}$ such that
	\[d'(x, c) < \delta' \implies d_{\RR}(f(x), f(c)) = |f(x) - f(c)| < \eps.\]
	
	\paragraph{Convergence.}
	Supposing that a sequence $(x_n)$ of points in $X$
	converges to some $\ell$ with respect to $d$, we wish to show that
	$(x_n)$ converges to $\ell$ with respect to $d'$ as well,
	meaning that for any $\eps > 0$,
	we want to show the existence of an $N$ such that
	\[n > N \implies d'(x_n, \ell) < \eps.\]
	
	By definition of convergence (with respect to $d$),
	there exists an $N_0$ such that
	\[n > N_0 \implies d(x_n, \ell) < a\eps,\]
	since $a\eps$ is a positive real number.
	In particular $n > N_0$ implies that
	\[a\cdot d'(x_n, \ell) \le d(x_n, \ell) < a\eps \implies
	d'(x_n, \ell) < \eps,\]
	so taking $N = N_0$ suffices to show that
	$(x_n)$ converges with respect to $d'$, too.
\end{soln}

\begin{problem}{5}
	Let $X = \RR^n$. For $\mathbf x = (x_1, \dots, x_n)$
	and $\mathbf y = (y_1, \dots, y_n)$ in $X$ set
	\[d_1(\mathbf x, \mathbf y) = \sum_{i=1}^n |x_i - y_i|
	\quad\text{and}\quad
	d_2(\mathbf x, \mathbf y) = \max_{i=1, \dots, n}|x_i - y_i|.\]
	Show that $d_1$ and $d_2$ are both distance functions on $X$.
	Show that the standard metric on $X$ we defined during class
	\[d(\mathbf x,\mathbf y) = \sqrt{\sum_{i=1}^n |x_i - y_i|^2}\]
	is equivalent to both $d_1$ and $d_2$.
\end{problem}

\begin{soln}
	First, we show that $d_1$ and $d_2$ are metrics on $X$.
	Indeed, both are clearly functions $X\times X\to\RR_{\ge0}$
	where $\RR_{\ge0}$ is the set of nonnegative reals
	(because absolute values and sums of absolute values
	are always nonnegative).
	We verify the other properties of metrics individually.
	First, for $d_1$:
	\begin{itemize}
		\ii $d_1(\mathbf x,\mathbf y) = 0$ iff $\mathbf x = \mathbf y$.
		First, clearly if $\mathbf x = \mathbf y$
		then $|x_i - y_i| = 0$ for all $i$, hence their sum
		$d_1(\mathbf x, \mathbf y)$ is also zero.
		
		Now, we prove the converse statement that
		if $\mathbf x\ne \mathbf y$, then the distance between them is not zero.
		To show this, we observe that $|x_i - y_i| \ge 0$ for all $i$,
		so, if, for some $k$,
		$|x_k - y_k|$ is positive (which $x_k\ne y_k$),
		then we are guaranteed to have
		\[\sum_{i=1}^n |x_i - y_i| \ge |x_k - y_k| > 0,\]
		meaning $d_1(\mathbf x,\mathbf y) \ne 0$.
		
		\ii $d_1(\mathbf x, \mathbf y) = d_1(\mathbf y, \mathbf x)$.
		This is fairly simple, as, because $|x_i - y_i| = |y_i - x_i|$, we have
		\[d_1(\mathbf x,\mathbf y) = \sum_{i=1}^n |x_i - y_i|
		= \sum_{i=1}^n |y_i - x_i| = d_1(\mathbf y, \mathbf x).\]
		
		\ii Triangle inequality. Let $\mathbf z = (z_1, \dots, z_n)$.
		Using the standard triangle inequality $|a - b| + |b - c| \ge |a - c|$,
		\[d_1(\mathbf x,\mathbf z) + d_1(\mathbf z,\mathbf y) =
		\sum_{i=1}^n |x_i - z_i| + |z_i - y_i| \ge
		\sum_{i=1}^n |x_i - y_i| = d_1(\mathbf x, \mathbf y).\]
	\end{itemize}
	For $d_2$:
	\begin{itemize}
		\ii $d_2(\mathbf x,\mathbf y) = 0$ iff $\mathbf x = \mathbf y$.
		First, clearly if $\mathbf x = \mathbf y$
		then $|x_i - y_i| = 0$ for all $i$, hence their max
		$d_2(\mathbf x, \mathbf y)$ is also zero.
		
		Now, we prove the converse statement that
		if $\mathbf x\ne \mathbf y$, then the distance between them is not zero.
		To show this, we observe that $|x_i - y_i| \ge 0$ for all $i$,
		so, if, for some $k$,
		$|x_k - y_k|$ is positive (which $x_k\ne y_k$),
		then we are guaranteed to have
		\[\max_{i=1,\dots,n} |x_i - y_i| \ge |x_k - y_k| > 0,\]
		meaning $d_2(\mathbf x,\mathbf y) \ne 0$.
		
		\ii $d_1(\mathbf x, \mathbf y) = d_1(\mathbf y, \mathbf x)$.
		This is fairly simple, as, because $|x_i - y_i| = |y_i - x_i|$, we have
		\[d_2(\mathbf x,\mathbf y) = \max_{i=1,\dots,n} |x_i - y_i|
		= \max_{i=1,\dots,n} |y_i - x_i| = d_2(\mathbf y, \mathbf x).\]
		
		\ii Triangle inequality. Let $\mathbf z = (z_1, \dots, z_n)$.
		Using the standard triangle inequality $|a - b| + |b - c| \ge |a - c|$,
		\begin{align*}
			d_2(\mathbf x,\mathbf y) &= \max_{i=1,\dots,n}|x_i - y_i|\\
			&\le \max_{i=1,\dots,n} (|x_i - z_i| + |z_i - y_i|)\\
			&\le \max_{i=1,\dots,n} \left(|x_i - z_i| + \max_{i=1,\dots,n}|z_i - y_i|\right)\\
			&= \max_{i=1,\dots,n} |x_i - z_i|
			+ \max_{i=1,\dots,n}|z_i - y_i|\\
			&= d_2(\mathbf x,\mathbf z) + d_2(\mathbf z,\mathbf y).
		\end{align*}
		To be clear, the second line is at most as large as the third line
		because the maximum of $|z_i - y_i|$ across all $i$
		is always at least as large as any one of the $|z_i - y_i|$,
		and the third line is equal to the fourth line
		because we are just separating the constant $\max |z_i - y_i|$
		from the outer max.
	\end{itemize}

	It remains to show that $d_1$ and $d_2$ are equivalent to $d$.
	We prove the following inequality,
	which clearly shows that the metrics are all equivalent.
	\[
	\frac{d_1(\mathbf x, \mathbf y)}{n} \le
	d_2(\mathbf x, \mathbf y) \le
	d(\mathbf x, \mathbf y) \le
	\sqrt n\cdot d_2(\mathbf x, \mathbf y) \le
	\sqrt n\cdot d_1(\mathbf x, \mathbf y).
	\]
	This splits into four inequalities in the obvious way;
	we prove them individually. First,
	\[\frac{d_1(\mathbf x, \mathbf y)}{n} =
	\frac{\sum_{i=1}^n |x_i - y_i|}{n} \le
	\frac{\sum_{i=1}^n\max_{i=1,\dots,n} |x_i - y_i|}{n}
	= \frac{n\cdot d_2(\mathbf x, \mathbf y)}{n}
	= d_2(\mathbf x, \mathbf y).\]
	Second,
	\[d_2(\mathbf x, \mathbf y) = \sqrt{\max_{i=1,\dots,n} |x_i - y_i|^2}
	\le \sqrt{\sum_{i=1}^n |x_i - y_i|^2} = d(\mathbf x, \mathbf y).\]
	Third,
	\[d(\mathbf x, \mathbf y) = \sqrt{\sum_{i=1}^n |x_i - y_i|^2} \le
	\sqrt{n\cdot \max_{i=1,\dots,n} |x_i - y_i|^2}
	= \sqrt n \cdot d_2(\mathbf x, \mathbf y).\]
	Fourth,
	\[\sqrt n\cdot d_2(\mathbf x, \mathbf y) =
	\sqrt n \cdot \max_{i=1,\dots,n}|x_i - y_i| \le
	\sqrt n \cdot \sum_{i=1}^n |x_i - y_i| =
	\sqrt n \cdot d_1(\mathbf x, \mathbf y).\]
	In particular we have that
	\[\frac{d_1(\mathbf x, \mathbf y)}{n} \le d(\mathbf x, \mathbf y) \le
	\sqrt n\cdot d_1(\mathbf x, \mathbf y),\]
	which shows that $d$ and $d_1$ are equivalent, and
	\[d_2(\mathbf x, \mathbf y) \le d(\mathbf x, \mathbf y) \le
	\sqrt n\cdot d_2(\mathbf x, \mathbf y),\]
	which shows that $d$ and $d_2$ are equivalent, as desired.
\end{soln}

\begin{problem}{6}
	Prove that if $(X, d_X)$ and $(Y, d_Y)$ are metric spaces
	then their Cartesian product $X\times Y$ is also a metric space
	when equipped with the metric
	\[d_{X\times Y} ((x, y), (x', y')) = d_X(x, x') + d_Y(y, y').\]
\end{problem}

\begin{soln}
	Since $d_X(x, x')$ and $d_Y(y, y')$ are always nonnegative,
	so is their sum; therefore $d_{X\times Y}$ is always nonnegative.
	Moreover,
	\[d_X(x, x') + d_Y(y, y') = 0\]
	if and only if $d_X(x, x')$ and $d_Y(y, y')$ are both zero,
	since if either is nonzero then the sum is positive;
	therefore, the sum $d_{X\times Y}((x, y), (x', y'))$ is zero
	if and only if $x = x'$ and $y = y'$.
	We also have
	\begin{align*}
		d_{X\times Y}((x, y), (x', y')) &=
		d_X(x, x') + d_Y(y, y')\\
		&= d_X(x', x) + d_Y(y', y)\\
		&= d_{X\times Y}((x', y'), (x, y)),
	\end{align*}
	and, for the triangle inequality,
	if we have a third point $(\hat x, \hat y)$, we can use the
	triangle inequalities for $d_X$ and $d_Y$ to get
	\begin{align*}
		d_{X\times Y}((x, y), (\hat x, \hat y)) +
		d_{X\times Y}((\hat x, \hat y), (x', y'))
		&= (d_X(x, \hat x) + d_Y(y, \hat y)) +
		(d_X(\hat x, x') + d_Y(\hat y, y'))\\
		&= (d_X(x, \hat x) + d_X(\hat x, x')) +
		(d_Y(y, \hat y) + d_Y(\hat y, y'))\\
		&\ge d_X(x, x') + d_Y(y, y')\\
		&= d_{X\times Y}((x, y), (x', y')).\qedhere
	\end{align*}
\end{soln}

\begin{problem}{7}
	Since $\RR^4 = \RR^3\times \RR = \RR^2\times\RR^2 = \RR\times\RR^3$
	there are three different product metrics on $\RR^4$.
	Write down an explicit formula for each one of those.
	What is the product metric on $\RR\times\RR\times\cdots\times\RR=\RR^n$?
\end{problem}

\begin{soln}
	Let $\mathbf x = (x_1, x_2, x_3, x_4)$ and
	$\mathbf y = (y_1, y_2, y_3, y_4)$.
	The product metric on $\RR^3\times\RR$ is
	\[d_{\RR^3\times\RR}(\mathbf x,\mathbf y) =
	\sqrt{|x_1 - y_1|^2 + |x_2 - y_2|^2 + |x_3 - y_3|^2}
	+ |x_4 - y_4|.\]
	The product metric on $\RR^2\times\RR^2$ is
	\[d_{\RR^2\times\RR^2}(\mathbf x,\mathbf y) =
	\sqrt{|x_1 - y_1|^2 + |x_2 - y_2|^2}
	+ \sqrt{|x_3 - y_3|^2 + |x_4 - y_4|^2}.\]
	The product metric on $\RR\times\RR^3$ is
	\[d_{\RR\times\RR^3}(\mathbf x,\mathbf y) =
	|x_1 - y_1| + \sqrt{|x_2 - y_2|^2 + |x_3 - y_3|^2 + |x_4 - y_4|^2}.\]
	For $\RR\times\RR\times\cdots\times\RR = \RR^n$,
	if $\mathbf x = (x_1, \dots, x_n)$ and $\mathbf y = (y_1, \dots, y_n)$, then
	\[d_{\RR\times\RR\times\cdots\times\RR} = \sum_{i=1}^n d_\RR(x_i, y_i)
	= \sum_{i=1}^n |x_i - y_i|.\qedhere\]
\end{soln}

\begin{problem}{8}
	Show that the product metric $d_{\RR^n\times\RR^m}$ on
	$\RR^n\times\RR^m = \RR^{n+m}$ is equivalent to the standard metric
	\[d(\mathbf x, \mathbf y) = \sqrt{\sum_{i=1}^{n+m} |x_i - y_i|^2}.\]
\end{problem}

\begin{soln}
	We write $d'$ instead of $d_{\RR^n\times\RR^m}$.
	We will prove the inequality
	\[\frac{1}{\sqrt 2} d'(\mathbf x, \mathbf y) \le
	d(\mathbf x, \mathbf y) \le d'(\mathbf x, \mathbf y),\]
	which by definition implies the metrics are equivalent. Define
	\[p = \sum_{i=1}^n |x_i - y_i|^2,\qquad q = \sum_{i=n+1}^{n+m} |x_i-y_i|^2\]
	so that
	\[d'(\mathbf x, \mathbf y) = \sqrt p + \sqrt q,\qquad
	d(\mathbf x, \mathbf y) = \sqrt{p+q}.\]
	Then, squaring gives the equivalent inequality
	\[\frac{p + q + 2\sqrt{pq}}{2} \le p + q \le p + q + 2\sqrt{pq}.\]
	The left inequality reduces to
	\[\sqrt{pq} \le \frac{p + q}{2}
	\iff 4pq \le p^2 + 2pq + q^2 \iff 0 \le p^2 - 2pq + q^2 = (p - q)^2\]
	which holds because all squares are nonnegative
	(or this is just the AM-GM inequality for two variables).
	The right inequality is trivial because $2\sqrt{pq} \ge 0$.
	Thus, the metrics are equivalent.
\end{soln}

\begin{problem}{9}
	In class we defined $\ell^2\subseteq\RR^\infty$ to be the set
	consisting of all infinite sequences $\mathbf x = (x_1, x_2, x_3, \dots)$
	such that
	\[\sum_{n=1}^\infty x_n^2 < +\infty.\]
	During class, for two sequences $\mathbf x$ and $\mathbf y$ in $\ell^2$
	we defined their distance to be
	\[d(\mathbf x, \mathbf y) = \sqrt{\sum_{n=1}^\infty |x_n - y_n|^2}.\]
	Explain why this series converges.
\end{problem}

\begin{soln}
	We want to show that
	\[\sum_{n=1}^\infty |x_n - y_n|^2 =
	\sum_{n=1}^\infty x_n^2 + y_n^2 - 2x_ny_n\]
	converges, given that $(x_n)$ and $(y_n)$ are sequences in $\ell^2$.
	
	\begin{claim*}
		The series
		\[\sum_{n=1}^\infty x_ny_n\]
		converges absolutely.
	\end{claim*}
	\begin{proof}
		Since $\mathbf x,\,\mathbf y\in\ell^2$, we know that
		\[\sum_{n=1}^\infty x_n^2,\qquad \sum_{n=1}^\infty y_n^2\]
		both converge.
		Moreover, we know that
		\[x_n^2 - 2|x_ny_n| + y_n^2 = (|x_n| - |y_n|)^2 \ge 0
		\iff x_n^2 + y_n^2 \ge 2|x_ny_n| \ge |x_ny_n|\]
		for all natural numbers $n$, so we can see by limit comparison that
		\[\sum_{n=1}^\infty |x_ny_n| \le \sum_{n=1}^\infty x_n^2 + y_n^2
		< +\infty,\]
		as desired.
	\end{proof}
	
	The claim in particular tells us that
	\[\sum_{n=1}^\infty x_ny_n\]
	converges (absolute convergence implies convergence), so we can say that
	\[\sum_{n=1}^\infty x_n^2 + y_n^2 - 2x_ny_n =
	\sum_{n=1}^\infty x_n^2 + \sum_{n=1}^\infty y_n^2 -
	2\sum_{n=1}^\infty x_ny_n < +\infty,\]
	hence the metric on $\ell^2$ is well-defined.
\end{soln}

\begin{problem}{10}
	Define $\ell^1\subseteq\RR^{\infty}$ to be the set of
	all infinite sequences $\mathbf x = (x_1, x_2, x_3, \dots)$
	that give rise to an absolutely convergent series.
	Show that $\ell^1$ is a metric space with distance
	\[d(\mathbf x, \mathbf y) = {\sum_{n=1}^{\infty} |x_n - y_n|}.\]
	Show that $\ell^1\subseteq\ell^2$.
	Find a sequence that is in $\ell^2$ but not in $\ell^1$
	(i.e. show that the inclusion is strict).
\end{problem}

\begin{soln}
	Clearly $d \ge 0$ over all of its inputs,
	and the distance function is symmetric because
	\[d(\mathbf x, \mathbf y) = \sum_{n=1}^\infty |x_n - y_n| =
	\sum_{n=1}^\infty |y_n - x_n| = d(\mathbf y, \mathbf x).\]
	We also see that $d(\mathbf x, \mathbf y) = 0$ if and only if
	\[\sum_{n=1}^\infty |x_n - y_n| = 0.\]
	Since each $|x_n - y_n|$ term is nonnegative,
	if any of the terms is positive then the sum will be positive,
	thus the only way the distance between two sequences is zero
	is if $|x_n - y_n| = 0$ for all $n$, i.e. $\mathbf x = \mathbf y$.
	
	To finish, the triangle inequality follows directly
	from the one-dimensional version:
	\begin{align*}
		d(\mathbf x, \mathbf z) + d(\mathbf z, \mathbf y) &=
		\sum_{n=1}^\infty |x_n - z_n| + \sum_{n=1}^\infty |z_n - y_n|\\
		&= \sum_{n=1}^\infty |x_n - z_n| + |z_n - y_n|\\
		&\ge \sum_{n=1}^\infty |x_n - y_n|\\
		&= d(\mathbf x, \mathbf y).
	\end{align*}
	Thus, we have shown that $d\colon\ell^1\times\ell^1\to[0,+\infty)$
	is zero if and only if its arguments are positive,
	is symmetric, and
	satisfies the triangle inequality,
	meaning it is a metric on $\ell^1$.
	
	Now, we show $\ell^1\subseteq\ell^2$.
	If $\mathbf x\in\ell^1$, then $|x_n| < 1$ eventually,
	and then $x_n^2 < |x_n| < 1$ eventually, say when $n > N$.
	By limit comparison,
	\[\sum_{n > N} x_n^2 < \sum_{n > N} |x_n| < +\infty,\]
	and $x_1^2$, \dots, $x_N^2$ don't affect the convergence
	of the series, so we can conclude that $\sum x_n^2$ converges.
	Additionally, the sequence given by $x_n = \frac 1n$
	is in $\ell^2$ but not $\ell^1$, as $\sum_{n=1}^\infty \frac 1n$
	diverges while $\sum_{n=1}^{\infty} \frac 1{n^2}$ converges.
\end{soln}


\end{document}